% Generated draft: requires mathematical and source review.
\documentclass[11pt]{article}
\usepackage[margin=25mm]{geometry}
\usepackage{fontspec}
\setmainfont{Latin Modern Roman}
\usepackage{amsmath,amssymb,mathtools,unicode-math}
\setmathfont{Latin Modern Math}
\usepackage{enumitem,needspace,xurl,hyperref}
\setlength{\parindent}{0pt}
\setlength{\parskip}{5pt}
\setlength{\emergencystretch}{3em}
\newcommand{\N}{\mathbb{N}}
\newcommand{\Z}{\mathbb{Z}}
\newcommand{\Q}{\mathbb{Q}}
\newcommand{\R}{\mathbb{R}}
\newcommand{\C}{\mathbb{C}}
\newcommand{\F}{\mathbb{F}}
\newcommand{\A}{\mathbb{A}}
\newcommand{\PP}{\mathbb{P}}
\newcommand{\E}{\mathbb{E}}
\newcommand{\eps}{\varepsilon}
\newcommand{\dd}{\,\mathrm{d}}
\begin{document}
\section*{UnsolvedMath ID: 700422}
\textbf{Source problem code:} catalogue entry
\subsection{Definitions and mathematical statement}
A group identity is a universally quantified equation between group words, equivalently an equation \(w(x_1,\ldots,x_n)=1\). For a list \(\Sigma\) of group identities, let
\[
\mathcal V(\Sigma)=\{G: G\text{ is a group satisfying every identity in }\Sigma\}.
\]
Here groups are arbitrary, with no finiteness assumption. In the first-order language of groups, write
\[
\operatorname{Th}(\mathcal V(\Sigma))
=\{\varphi:\varphi\text{ is a sentence and }G\models\varphi\text{ for every }G\in\mathcal V(\Sigma)\}.
\]
This theory is called decidable if an algorithm, given any sentence \(\varphi\), terminates and decides whether \(\varphi\) belongs to it. The source's word ``soluble'' is understood here as decidable, rather than as a group-theoretic solvability condition.

\textbf{Question (A. V. Kuznetsov, 1980).} Does there exist an algorithm which, given a list \(\Sigma\) of group identities, terminates and decides whether \(\operatorname{Th}(\mathcal V(\Sigma))\) is decidable?

The source attributes to A. P. Zamyatin the equivalence between this decidability condition and the condition that the variety is abelian. With that supplied equivalence, the question asks for an algorithm deciding whether
\[
\forall G\in\mathcal V(\Sigma)\;\forall x,y\in G,\qquad xy=yx.
\]
Equivalently, it asks whether the commutativity identity holds in every group satisfying \(\Sigma\); it does not require \(\mathcal V(\Sigma)\) to contain every abelian group.

\textbf{Input qualification.} A finite list of finitely written identities gives the usual finite algorithmic input. Under that interpretation the algorithm must work for every such list, including the empty list. The supplied wording does not explicitly restrict lists to be finite or specify an effective presentation for infinite lists. Thus the finite-list formulation is an explicitly stated interpretation; a formulation allowing infinite lists requires additional input conventions. The cited equivalence and its precise hypotheses have not been independently checked.
\subsection{Short English statement}
Kuznetsov asks whether an algorithm can determine, from defining group identities, whether the resulting variety has a decidable elementary theory. The source cites Zamyatin for an equivalent test: whether every group satisfying those identities is abelian. The finite-list input convention is implicit in this draft and requires review. The supplied registry labels the problem open; that label is not an independently verified assessment.
\subsection{Sources}
\begin{itemize}
\item Local UnsolvedMath registry, permanent ID \(700422\), exact source code \texttt{catalogue entry}, title \emph{Kourovka Notebook Problem 7.23}. Supplied upstream revision: \texttt{d371ee917ac9647fbb06cd6fc529deec6f3dda82}. The supplied record labels the problem \emph{open}.
\item \emph{The Kourovka Notebook (Unsolved Problems in Group Theory)}, original problem 7.23, 7th issue (1980); proposed by A. V. Kuznetsov. The supplied background identifies the 21st issue edition (2026), arXiv:1401.0300v48. The registry describes the collection as Issues 1--21 (1965--2026), edited by E. I. Khukhro and V. D. Mazurov.
\item A. P. Zamyatin, \emph{Algebra and Logic}, \textbf{17}, no.\ 1 (1978), 13--17, as cited in the supplied problem statement for the equivalence between decidability of the elementary theory and the variety being abelian. No article title or further bibliographic details are supplied. This attribution and its hypotheses have not been independently verified.
\end{itemize}
\end{document}
